\section{Matroid Products}\label{sec:MatProds}

In this section we will revisit some ideas of Mason \cite{MGlueing} and Las Vergnas \cite{LV81} particularly with respect to tensor and symmetric powers of matroids, working with matroids rather than valuated matroids.
   
Work with matroid products in the category of matroids with strong maps was explored by Crapo and Rota in \cite{CR70}, and it was known at the time of these works that matroid products did not exist in this category. However, this is not the notion of matroid product that we are discussing in this paper. Instead, we interest ourselves in matroid products in a weaker sense, as defined by Las Vergnas for tensor products, and by Mason for symmetric powers.
\subsection{Tensor Products}\label{subsec:tensorProducts}
    Given matroids $M$ and $N$, we aim to construct a new matroid, called a tensor product of $M$ and $N$, which will be a matroid with ground set equal to the Cartesian product $E(M)\times E(N)$ with rank $\rk(M)\rk(N)$ and inheriting some structure from $M$ and $N$. While this operation is always possible for matroids represented over a common field $k$, it is not necessarily unique, nor is it  possible for matroids in general as shown by LasVergnas in \cite{LV81}*{Proposition 2.1}. 
    
    Several results presented in this section have been known since LasVergnas's original paper, however, we revisit this work and introduce a new proof of \cite{LV81}*{Proposition 2.1}. Furthermore, our proof generalizes to an infinte class of non-algebraic matroids with no common non-realizable minors. This family of matroids is exactly the class of matroids that were initially studied by Lindstr\"om in \cite{Lind88}. We substantiate the claim that none of these matroids has a non-trivial matroid tensor product.
    
    \begin{definition}\label{def:tensorproduct}
    Let $M$ and $N$ be matroids on ground sets $E(M)$ and $E(N)$ respectively. A matroid $P$ on ground set $E(M)\times E(N)$ is called a \textbf{quasi product} of $M$ and $N$ if for every non-loop $e\in E(M)$ we have
    \begin{equation} \label{eqn:tensorRight}
        N\cong P|_{\set{(e,x)\mid x\in E(N)}}
    \end{equation}
    and for every non-loop $f\in E(N)$ we have
    \begin{equation}\label{eqn:tensorLeft}
        M\cong P|_{\set{(x,f)\mid x\in E(M)}},
    \end{equation}
    where these isomorphisms are explicitly given by the maps $x\mapsto (e,x)$ and $x\mapsto (x,f)$ respectively. If $e$ is a loop in $E(M)$ and $f$ a loop in $E(N)$ then we have
    \begin{equation}\label{eqn:tensorRight0}
        \rk(P|_{\set{(e,x)\mid x\in E(N)}})=0
    \end{equation}
    or
    \begin{equation}\label{eqn:tensorLeft0}
    \rk(P|_{\set{(x,f)\mid x\in E(M)}})=0
    \end{equation}
    respectively.

    We will refer to the set $\set{(e,x)\mid x\in E(N)}$ as the \textbf{row} associated to $e\in E(M)$ and we will refer to the set $\set{(x,f)\mid x\in E(M)}$ as the \textbf{column} associated to an element $f\in E(N)$ .
    \end{definition}
    
    \hyperref[eqn:tensorRight]{Isomorphism \ref*{eqn:tensorRight}} and \hyperref[eqn:tensorLeft]{Isomorphism \ref*{eqn:tensorLeft}} imply that, given bases $B_1$ of $M$ and $B_2$ of $N$, that $B_1\times B_2$ is a spanning set of any quasi product of $M$ and $N$. Therefore the rank of a quasi product is bounded above by $\rk(M)\rk(N)$, and this bound can be obtained whenever $M$ and $N$ are $k$-realizable matroids. We call that a quasi product of $M$ and $N$ with rank $\rk(M)\rk(N)$ a \textbf{tensor product} of $M$ and $N$. Moreover, if $M$ and $N$ are the same matroid, then we refer to any of their tensor products as a \textbf{second tensor power}.

       
    The motivating example for the definition of matroid tensor products is, as is often the case, linear spaces. Given two matroids $M_1$ and $M_2$ that are represented as linear subspaces of $L(M_1)\subset k^{d_1}$ and $L(M_2)\subset k^{d_2}$, the tensor product $L(M_1)\otimes L(M_2)$ is a linear subspace of $k^{d_1d_2}$ whose corresponding matroid is a tensor product of $M_1$ and $M_2$. Given matrix representations of $L(M_1)$ and $L(M_2)$ their Kronecker product yields a matrix representation of $L(M_1)\otimes L(M_2)$.
    
    The following lemma is a result of LasVergnas, and is useful for understanding the structure of tensor products.
    \begin{lemma}\label{lem:tensorCharacterizations}\cite{LV81}*{Lemma 2.2}
    Let $M$ and $N$ be matroids and let $P$ be a quasi product of $M$ and $N$. The following are equivalent.
    \begin{enumerate}
        \item $\rk(P)=\rk(M)\rk(N).$
        \item For every basis $B_M$ of $M$ and every basis $B_N$ of $N$, $B_M\times B_N$ is a basis of $P$.
        \item For every pair of flats $F'\subset F$ of $M$ and $G'\subset G$ of $N$ 
        $$F'\times G \cup F\times G'$$
        is a flat of $P$.
        \item For every pair of sets $X'\subset X\subset E_M$ and $Y'\subset Y\subset E_N$
        $$\rk(X'\times Y\cup X\times Y')= \rk(X'\times Y)+\rk(X\times Y)- \rk(X'\times Y').$$
    \end{enumerate}
    \end{lemma}
    
    It is important to keep in mind that matroid tensor products may not be unique. In fact, even for a fixed field $k$ and two $k$-representable matroids $M_1$ and $M_2$, the tensor product obtained from linear subspace representations $L(M_1)$ and $L(M_2)$ may depend on the choice of representatives $L(M_1)$ and $L(M_2)$. 

    However, existence and uniqueness are guaranteed in some cases.

    \begin{lemma}\label{lem:tensorSums}
     Let $A,B,$ and $M$ be any matroids and let $N=A\oplus B$. If $P$ is a tensor product of $M$ and $N$ then $P= P_A\oplus P_B$ where $P_A$ is a tensor product of $M$ and $A$ and $P_B$ is a tensor product of $M$ and $B$.
    \end{lemma}
    \begin{proof}
        Consider the matroids 
        $$P_A = P|_{E(M)\times E(A)},$$
        and
        $$P_B = P|_{E(M)\times E(A)}.$$
        It is straightforward to check that $P_A$ and $P_B$ are quasi products of $M$ with $A$ and $B$ respectively. Since
        $$\rk(P)=\rank(M)\rank(A)+\rank(M)\rank(B))$$
        we may conclude that $\rk(P_A)=\rk(M)\rk(A)$ and $\rk(P_B)= \rk(M)\rk(B)$. Therefore $P_A$ and $P_B$ are tensor products. 
    \end{proof}

    \begin{lemma}\label{lem:tensorSimplification}
    Let $M$ and $N$ be any matroids. Then any quasi product of $M$ with $N$ depends only the simplification of $M$. 
    \end{lemma}
    \begin{proof}
        Let $P$ be a quasi product of $M$ with $N$. Suppose $M$ has a loop $e$. Then for any $x$ in $E(N)$, we have that $(e,x)$ is a loop of $P$. This is an immediate consequence of \hyperref[eqn:tensorRight0]{Equation \ref*{eqn:tensorRight0}}. In particular, adding or removing a loop $e$ from $M$ affects the quasi product $P$ by adding or removing the set of loops $\set{e}\times E(N)$ from $P$. 

        Now suppose that $S$ is a parallelism class of $M$. Then \hyperref[eqn:tensorRight]{Isomorphism \ref*{eqn:tensorRight}} implies that $S\times \set{x}$ is contained in a parallelism class $S'$ of $P$. In particular, replacing the parallelism class $S$ with a representative $e$ affects $P$ by replacing the subset $S\times \set{x}$ of the parallelism class $S'$ with $(e,x)$. 
    \end{proof}
\begin{proposition}\label{prop:uniqueTensors}
Let $M$ be any matroid. Then there is a unique tensor product between $M$ and $U_{n,n}$.
\end{proposition}
    \begin{proof}
            If $P$ is a tensor product of $M$ and $U_{n,n}$, then \hyperref[lem:tensorSums]{Proposition \ref*{lem:tensorSums}} implies that $P = \oplus_{i=1}^n P_i$ where $P_i$ is a tensor product of $M$ and $U_{1,1}$. In particular $P_i\cong M$ and so $P=\oplus_{i=1}^n M$. We see that $\rk(P)=n\rk(M)=\rk(U_{n,n})\rk(M)$ and so $P$ is tensor product. This proves existence and uniqueness.
    \end{proof}

    \begin{corollary}\label{cor:uniqueSimplification}
        Let $M$ be any matroid and let $N$ be a matroid whose simplification is $U_{n,n}$. Then a tensor product between $M$ and $N$ exists and is unique.
    \end{corollary}
    \begin{proof}
    This follows from Propositions \hyperref[lem:tensorSums]{Proposition \ref*{lem:tensorSums}}, \hyperref[prop:uniqueTensors]{Proposition \ref*{prop:uniqueTensors}}, and \hyperref[lem:tensorSimplification]{Lemma \ref*{lem:tensorSimplification}}
    \end{proof}
 
    Outside of the realizable case, little is known about the existence of matroid tensor products beyond \hyperref[cor:uniqueSimplification]{Corollary \ref*{cor:uniqueSimplification}}. It is known, however, that not all matroids have tensor products between them. 
    
    In his first work on matroid tensor products, Las Vergnas proved that the V\'amos matroid $V_8$ has no second tensor power. However, in \cite{LV81} this proof is simplified into two smaller results. First, \cite{LV81}*{Proposition 2.1} proves that no tensor product between $U_{2,3}$ and $V_8$ exists. Then \cite{LV81}*{Remark 2.3} provides that if a tensor power $P$ existed, it would have a minor isomorphic to a tensor product of $U_{2,3}$ and $V_8$, a contradiction. We conclude that no second tensor power of $V_8$ exists. In particular, the V\'amos matroid is a forbidden minor of the class of matroids with second tensor powers. We make this connection between tensor products and their minors explicit.

    We denote the collection of independent sets of a matroid $N$ by $\mathbf{\sI(N)}$.

    \begin{proposition}\label{prop:tensorMinorClosed}
        Let $M$ and $N$ be any matroids and let $P$ be a tensor product of $M$ and $N$. Let $S$ and $T$ be disjoint subsets of $E(M)$ and let $M'= M\sm S\contr T$ be a minor of $M$. Then the minor $P':=P\sm (S\times E(N))\contr (T\times E(N))$ is a tensor product of $M'$ with $N$.
    \end{proposition}
    \begin{proof}
        It suffices to prove the result for single-element deletions and contractions. Moreover, by \hyperref[lem:tensorSimplification]{Lemma \ref{lem:tensorSimplification}}, we may assume that $M$, $N$, and $M\sm e$ are simple matroids.

        First consider the minor $P':=P\sm(e\times E(N))$. We prove that $P'$ is a tensor product of $M\sm e$ and $N$.
        
        To see that $P'$ is a quasi product, first note that for $a\in E(M)\sm e$ we have $P'|_a\cong P|_a\cong N$, so \hyperref[eqn:tensorRight]{Isomorphism \ref*{eqn:tensorRight}} holds. Furthermore, for $b\in E(N)$ we have $P'|_b\cong (P|_b)\sm (e\times b)\cong M\sm e$, as desired. Now let $I_M$ be independent in $M\sm e$ and $I_N$ independent in $N$. Then $I_M$ is independent in $M$, and therefore $I_M\times I_N$ is independent in $P$. This implies $I_M\times I_N$ is independent int $P'$. We conclude by \hyperref[lem:tensorCharacterizations]{Lemma \ref*{lem:tensorCharacterizations}} that $P'$ is a tensor product of $M\sm e$ and $N$.

        Now let $P':=P\contr (e\times E(N))$. We prove that $P'$ is a tensor product of $M\contr e$ and $N$. Consider the collection 
        $$a\times \sI(N):=\set{a\times I: I\in \sI(N)}.$$
        We prove that $\sI(P'|_a)= a\times \sI(N)$. First let $T$ be independent in $P'|_a$. Then there is an independent set $J'\in \sI(P')$ with $J'|_a=T$. Moreover, this implies there is an independent subset $J\in \sI(P)$ contained in $J'\cup (e\times E(N))$. In particular, $J|_a=T$ is an independent set in $P|_a\cong N$, so $T= a\times I$ for some $I\in \sI(N)$.

        For the reverse inclusion, suppose that $I$ is independent in $\sI(N)$ and let $B$ be a basis containing $I$. Then the set $(a\times I)\cup (e\times B)$ is contained in $\set{a,e}\times B$, which is independent in $P$. In particular, this implies $a\times I$ is independent in $P'$ and thus independent in $P'|_a$. This implies that \hyperref[eqn:tensorRight]{Isomorphism \ref*{eqn:tensorRight}} holds.

        Now let $b\in E(N)$. We proceed as before. First let $I_M$ be independent in $M\sm e$, and let $B_N$ be a basis of $N$ that contains $b$. Then $(I_M\cup e)\times B$ is independent in $P$ and contains the set $(I_M\times b)\cup (e\times B_N)$. In particular, $(I_M\times b)$ is independent in $P'$, and therefore in $P'|_b$.

        
        Now let $T$ be independent in $P'|_b$, say $T=S\times b$ for $S\subset E(M)\sm e$. Consider that since $T$ is independent in $P'|_b$, it is independent in $P'$, and therefore for any basis $B_N$ of $N$, the set $T\cup (e\times B_N)$ is independent in $P$. Suppose then that $B_N$ contains $b$. Then $\left(T\cup(e\times B_N)\right)_b = T\cup e$ is an independent set in $P|_b$. In particular $S\cup e$ is independent in $M$ and therefore $S$ is independent in $M\contr e$, as desired. We conclude $P'$ satisfies \hyperref[eqn:tensorLeft]{Isomorphism \ref{eqn:tensorLeft}}.

        We now check that $P'$ is indeed a tensor product. Let $I_M$ be independent in $M\contr e$ and $I_N$ independent in $N$. We claim that $I_M\times I_N$ is independent in $P'$. To see this, let $B_N$ be a basis of $N$ containing $I_N$, then $(I_M\cup e)\times B_N$ is independent in $P$, or in particular, $(I_M\times I_N) \cup (e\times B_n)$ is independent in $P$. This implies $I_M\times I_N$ is independent in $P'$. \hyperref[lem:tensorCharacterizations]{Lemma \ref{lem:tensorCharacterizations}} now implies that $P'$ is a tensor product of $M\contr e$ and $N$, as desired.
    \end{proof}

    We now provide a generalization of the proof that the V\'amos matroid has no second tensor power, and we further extend this result to the infinite family of non-algebraic matroids introduced by Lindstr\"om \cite{Lind88}. The matroid $M(r)$ is a rank $r$ matroid defined as follows. Let $A:=\set{a_1,\dots, a_{r-2}}$, $B:=\set{b_1,\dots, b_{r-2}}$, $C:=\set{c_1,\dots, c_{r-2}}$, and $D:=\set{d_1,\dots, d_{r-2}}$ be disjoint sets of size $r-2$. We let the ground set of $M(r)$ be $E_r:= A\cup B\cup C\cup D$. We define the set of \textbf{cyclic flats} of $M(r)$ --- the flats that are unions of circuits --- as $\sH:= \set{A\cup B, A\cup C, A\cup D, B\cup C, B\cup D}$. The cyclic flats of $M(r)$, together with their rank, determine $M$ as a matroid \cite{Bry75}. In the case of $M(r)$, every cyclic flat in $\sH$ has rank $r-1$ and so they are hyperplanes. Observe that $M(4)$ is the V\'amos matroid.

    \begin{proposition}\label{prop:XandY}
        Let $[3]$ denote the ground set of $U_{2,3}$ and let $E:=A\cup B\cup C\cup D$ denote the ground set of $M(r)$. Let $P$ be any quasi product of $U_{2,3}$ and $M(r)$. Then either the set
        $$X:=([3]\times A)\cup (1\times B)\cup (2\times C)\cup (3\times D)$$
        or the set
        $$Y:= ([3]\times B)\cup (1\times A)\cup (2\times C)\cup (3\times D),$$
        is a spanning set.
    \end{proposition}
    \begin{proof}
    
        Let $P$ be a quasi product of  $U_{2,3}$ and $M(r)$. If one of $X$ or $Y$ is a spanning set, we're done. So suppose neither $X$ nor $Y$ is spanning. We claim that $X$ and $Y$ are both hyperplanes of $P$. To see this, we show that adding any element to either $X$ or $Y$ results in a spanning set. To make this more clear we refer to the ground set in terms of rows, indexed by the ground set $[3]$ of $U_{2,3}$, and columns indexed by the ground set $E$ of $M(r)$. We first consider $X$. Adding an element $e=(e_1,e_2)$ to $X$ leads to a subset of the form
            $$e_1\times (A\cup T\cup e_2),$$
    where $T\in \set{B, C, D}$. Now the closure of $X\cup e$ necessarily contains all of row $e_1$ as $A\cup T\cup e_2$ is a spanning set of $M(r)$. This now implies that $\ol{X\cup e}$ contains at least $2$ elements in each column of $[3]\times E$, which provides a spanning set of $U_{2,3}$ in each column. Thus the closure in every column is the entire column and so $\ol{X\cup e}=[3]\times E$. A similar argument holds for $Y$, where $Y\cup e$ contains the row $e_1\times (B\cup T\cup e_2)$ where $T\in{A,C,D}$. We conclude that $X$ and $Y$ are hyperplanes of $P$. 

    
    Now consider that $Z:= X\cap Y$ is also a corank-1 set. This follows as as adding the element $e = (1,c_1)$ leads to a spanning set in row 1. Following the previous argument we see that $\ol{Z\cup (1,c_1)}$ contains at least 2 elements in the columns of $C$ and $D$, which is a spanning set in those column. So $\ol{(Z\cup(1,c_1))}$ contains the set $[3]\times (C\cup D)$. However, $C\cup D$ is a spanning set of $M(r)$, and so every row of $\ol{(Z\cup(1,c_1))}$ contains a spanning set. Therefore $\ol{(Z\cup(1,c_1))}= [3]\times E$. This is a contradiction as the intersection of two distinct hyperplanes, namely $X$ and $Y$, must have corank at least $2$, whereas $Z$ has corank $1$. Thus either $X$ or $Y$ must be a spanning set.
    \end{proof}
    
    \begin{proposition}\label{prop:VxU23}
    Any quasi product of $U_{2,3}$ and $M(r)$ has rank at most $2r-1$. In particular, no tensor product between $U_{2,3}$ and $M(r)$ exists.
    \end{proposition}
    \begin{proof}
        Let $[3]$ denote the ground set of $U_{2,3}$ and let $E:=A\cup B\cup C\cup D$ denote the ground set of $M(r)$. Let $P$ be a quasi product of $U_{2,3}$ and $M(r)$. Suppose $X$ and $Y$ are as in \hyperref[prop:XandY]{Proposition \ref*{prop:XandY}}, and without loss of generatlity let $X$ be a spanning set among $X$ and $Y$. We prove $X$ has rank at most $2r-1$. To see this, consider that $X$ is necessarily spanned by $([2]\times A)\cup \set{(1,b), (2,c), (3,d)}$, which contains $2r-1$ elements, thus $\rk(P)=\rk(X)\leq 2r-1$.
    \end{proof}
    
    As a consequence of \hyperref[prop:VxU23]{Proposition \ref*{prop:VxU23}}, the matroids $M(r)$ have no tensor powers, and we further note the following small corollary
    
    \begin{corollary}\label{cor:VxU23}
    A tensor product between $M(r)$ and $N$ exists if and only if the simplification of $N$ is a free matroid.
    \end{corollary}
    \begin{proof}
    If the simplification of $N$ is not free, then $N$ contains a circuit of size at least 3, and therefore $N$ contains $U_{2,3}$ as a minor. By \hyperref[prop:tensorMinorClosed]{Proposition \ref*{prop:tensorMinorClosed}} and \hyperref[prop:VxU23]{Proposition \ref*{prop:VxU23}}, no tensor product between $M(r)$ and $N$ can exist. Otherwise, the simplification of $N$ is free and \hyperref[cor:uniqueSimplification]{Corollary \ref*{cor:uniqueSimplification}} implies that $M(r)$ and $N$ have a unique tensor product.
    \end{proof}
    
    In addition to showing that each $M(r)$ has no tensor product with $U_{2,3}$ we draw the further conclusion that there is an infinite family of forbidden minors to the class of matroids with second tensor power, which is minor closed by \hyperref[prop:tensorMinorClosed]{Proposition \ref*{prop:tensorMinorClosed}}.
    \begin{corollary}\label{cor:tensorForbiddenMinors}
        The class of matroids with second tensor power has infinitely many forbidden minors.
    \end{corollary}
    \begin{proof}
    The matroids $M(r)$ satisfy the following properties \cite{Lind88}
    \begin{enumerate}
        \item $M(r)$ is a non-algebraic matroid.
        \item For any $e\in E_r$, the contraction $M(r) \contr e$ is realizable.
        \item If $r_1\neq r_2$, and two restrictions satisfy $M(r_1)|_{T_1}\cong M(r_2)|_{T_2}$ then $M(r_1)|_{T_1}\cong U_{n,n}$ for $n= |T_1|$.
    \end{enumerate}
    We deduce from these observations that, for $r_1\neq r_2$, any minor common between $M(r_1)$ and $M(r_2)$ is a realizable matroid. In particular, any common minor has second tensor powers. We conclude that each $M(r)$ has a minor, which is not a minor to any other $M(r')$ that is a forbidden minor of the class of matroids with second tensor powers.
    \end{proof}

    While the connection between tensor products and tropical geometry is not as immediate as will be seen with symmetric powers, we are nonetheless able to use tensor products in order to prove certain symmetric powers do not exist, as done in \cite{DR19}.

\subsection{Matroid Symmetric Powers}
We now turn our focus onto symmetric powers of matroids. Recall that, for any set $E$, the $d\th$ symmetric power of $E$ is the set of unordered $d$-tuples of elements of $E$, which we denote $\Sym_d(E)$. If we fix an element $\alpha \in \Sym_i(E)$, there is an injective ``multiplication" map
    \begin{align*}
        *_\alpha: \Sym_d(E)&\to \Sym_{d+i}(E)\\
        \beta&\mapsto (\alpha,\beta).
    \end{align*}
Since the coordinates of elements of $\Sym_{d+i}$ are unordered, we will write $(\alpha,\beta)$ as $\alpha\beta$, taking after commutative multiplication notation. 

We will often consider the subset of elements of $\Sym_{d+i}(E)$ that are divisible by some $\alpha\in \Sym_{i}(E)$, which we denote as

$$\Sym_{d+i}(E)|_{\alpha} := \set{e\in \Sym_{d+i}(E)\mid e= \alpha x}.$$

Given a valuated matroid $\sV$ on $\Sym_{d+i}(E)$ we further extend this notation and write $\sV|_{\alpha}$ instead of $\sV|_{\Sym_{d+1}(E)|_{\alpha}},$ where the restriction of $\sV$ to the set $T\subset [n]$ is given by the subsemimodule 
$$S|_T:=\sV\cap \set{x\in \T^n \mid x_i=\infty \text{ for all } i\notin T}.$$

\begin{definition}\label{def:symprod}
Fixing a matroid $\sV$ on a ground set $E$, a $\mathbf{d}^{\text{\textbf{th}}}$ \textbf{symmetric quasi power} of $\sV$ is a sequence of matroids $(\sV_1,\dots,\sV_d)$ with $\sV_1 = \sV$ and for all $1\leq i \leq d$
\begin{itemize}
    \item $E(\sV_i) = \Sym_i(E),$
    \item for all $\alpha\in \Sym_{d-i}(E)$ for which $\alpha$ is not divisible by a loop the multiplication map $*_\alpha$ given by $y\mapsto \alpha x$ induces the isomorphism
    \begin{equation}\label{eqn:symprodEmbedding} 
        \sV_i\cong \sV_d|_{\alpha}
    \end{equation}
    otherwise
    \begin{equation}\label{eqn:symprodLoops}
        \rk(\sV_d|_{\alpha})=0.
    \end{equation}
    
\end{itemize} 
Moreover, if $(\sV_1,\dots,\sV_d)$ is a symmetric quasi power of $\sV$ with

\begin{equation}\label{eqn:symprodRank}
\rk(\sV_d)= \binom{\rk(\sV)+d-1}{d}
\end{equation}

then we say that $(\sV_1,\dots,\sV_d)$ is a $\mathbf{d}^{\text{\textbf{th}}}$ \textbf{symmetric power} of $\sV$.
\end{definition}

The most familiar and pertinent example of symmetric powers arises from commutative polynomial rings. Given a polynomial ring $k[x_1,\dots,x_n]$ over a field (or semifield) $k$, the degree $d$ monomials in $k[x_1,\dots,x_n]$ exactly correspond to the set $\Sym_d(\set{x_1,\dots,x_n}),$ which we denote by $\Mon_d(n)$ --- the degree $d$ monomials in $n$ variables. The multiplication map $*_\alpha$ is then given by multiplication by the monomial $x^\alpha$ in $k[x_1,\dots,x_n]$.

In the setting of tropical polynomial semirings and symmetric powers, matroids can be constructed as collections of linear tropical polynomials, and their $d\th$ symmetric powers as collections of degree-$d$ tropical polynomials.

\begin{remark}
In \coloref{def:symprod} we define a $d\th$ symmetric power of a fixed valuated matroid $\sV$ to be a sequence $(\sV_1,\dots,\sV_d)$. However, \coloref{eqn:symprodEmbedding} provides that $\sV_d$ uniquely determines the entire sequence $(\sV_1,\dots,\sV_d)$, so nothing is lost when referring to $\sV_d$ alone. However, $\sV$ may have multiple $d\th$ symmetric powers, as we will see in \coloref{ex:matrixSymprodExample}, and when constructing a $d\th$ symmetric power it will be useful to explicitly record the entire sequence $(\sV_1,\dots,\sV_d)$. We will further discuss the space of all symmetric powers of a fixed valuated matroid $\sV$ in \coloref{subsec:SpacesofSymPow}.
\end{remark}

\begin{remark}\label{rem:truncation}
If $(\sV_1,\dots,\sV_d)$ is a $d\th$ symmetric power of $\sV$, then for every $1\leq i\leq d$ we have that the truncation $(\sV_1,\dots,\sV_i)$ is an $i\th$ symmetric power of $\sV$.
\end{remark}



\begin{remark}\label{rem:loopless}
\coloref{eqn:symprodEmbedding} implies that if $\sV$ is loop-free, then so is every symmetric quasi power of $\sV$. Moreover, \coloref{eqn:symprodLoops} implies that adding or removing loops to $\sV$ to obtain a matroid $\sV'$ uniquely determines how to add or remove loops from any given $d\th$ symmetric quasi power $(\sV_1,\dots,\sV_d)$ to obtain a $d\th$ symmetric quasi power of $\sV'$. In particular, we may assume without loss of generality that $\sV$ is loop-free and all subsequent results extend to matroids with loops.
\end{remark}

It is also important to note that we refer to \textit{a} symmetric power of a matroid rather than \textit{the} symmetric power, as it may not be unique, not even for realizable matroids.

The following is stated in \cite{MGlueing}, but we restate it here and include a brief proof for intuition's sake.

\begin{proposition}\cite{MGlueing}\label{prop:masonFullRank}
    For a matroid $M$ of rank $r$ the rank of any symmetric quasi power $M_d$ is at most $\binom{r+d-1}{d}$ and at least $r$.
\end{proposition}

\begin{proof}
For any symmetric quasi power $M_d$ of $M$ and any basis $B$ of $M$, \coloref{eqn:symprodEmbedding} implies that the set $\Sym_d(B)$ is a spanning set of $M_d$. Thus the rank of $M_d$ can never be greater than $|\Sym_d(B)|=\binom{r+d-1}{d}$.

\coloref{eqn:symprodEmbedding} further implies that for any $e\in E(M)$ and any basis $B\in \sB(M)$, that $eB$ is independent in $M_d$. In particular, $\rk(M_d)\geq r$.
\end{proof}
 In light of \coloref{prop:masonFullRank}, we see that symmetric powers are exactly those symmetric quasi-powers that achieve the upper bound on the rank. Moreover given any field $k$, every $k$-representable matroid $M$ has a $k$-representable $d\th$ symmetric power that can be directly constructed. Given any $k$-matrix $L$ representing $M$ we can construct the $d\th$ \textbf{Macaulay matrix} of $L$, which will be a representation of a $d\th$ symmetric power of $M$. We outline the construction of the second Macaulay matrix here. Let $r\leq n$ and let $L$ be a rank-$r$ $r\times n$ matrix over a field $k$. The second Macaulay matrix of $L$, which we will denote $\Sym_2(L)$, is an $\binom{r+1}{2}\times \binom{n+1}{2}$ matrix whose rows are indexed by $\Sym_2([r])$ and whose columns are indexed by $\Sym_2([n])$. Fixing the standard basis $e_1,\dots,e_r$ of $k^r$, we can write each column $C_j$ of $L$ as a formal sum 
$$C_j = \sum_{i=1}^r L_{i,j}e_i.$$
We can now consider $C_j$ as a linear form in the polynomial ring $k[u_1,\dots,u_r]$ under the inclusion $\phi:k^n\mapsto k[u_1,\dots,u_r]$, $e_i\mapsto u_i$.

The column $ij$ of $\Sym_2(L)$ is computed by taking the product $\phi(C_i)\phi(C_j)\in k[u_1,\dots,u_r]_2$, which is then a formal sum in the basis $\Mon_2(r)$, i.e in basis vectors indexed by $\Sym_2([r])$. In particular we express 
$$\phi(C_i)\phi(C_j)=\sum_{x,y\in[r]} c_{xy,ij}u_xu_y.$$

We now define the entries of the second Macaulay matrix $\Sym_2(L)$ to be, for $x,y\in [r]$ and $i,j\in[n]$
$$\Sym_2(L)_{xy,ij} = c_{xy,ij},$$
that is, the $xy$ coefficient of the product of the polynomials associated to columns $C_i$ and $C_j$. This definition extends to defining $\Sym_{d}(L)$ by considering products in $k[x_1,\dots,x_n]_d.$

As indicated in our notation above, the $d\th$ Macaulay matrix of $L$ is indeed a matrix that represents a $d\th$ symmetric power of the matroid $M$ that $L$ represents. We will demonstrate this for the second Macaulay matrix. 
\begin{example}\label{ex:matrixSymprodExample}
Consider the matrix representing $U_{2,4}$ over any field other than $k_2$ with $\alpha \neq 0,1$
$$M:=\begin{bmatrix}
1 & 0 & 1 & 1\\
0 & 1 & 1 & \alpha
\end{bmatrix}.$$
The second Macaulay matrix of $M$ is a matrix whose column space represents a second symmetric power of $M$. It's columns are indexed by pairs of columns of $M$ and its rows are indexed by pairs of column entries of some fixed basis of $M$, in this case we use the first two columns of $M$.
$$\Sym_2(M):=\kbordermatrix{
        & c_{11}& c_{12}& c_{13}& c_{14} & c_{22}& c_{23}& c_{24} & c_{33}& c_{34}  & c_{44}  \\
r_{11}  &1      & 0     & 1     & 1      & 0     & 0     & 0      & 1     & 1       & 1       \\
r_{12}  &0      & 1     & 1     & \alpha & 0     & 1     & 1      & 2     & \alpha+1& 2\alpha \\
r_{22}  &0      & 0     & 0     & 0      & 1     & 1     & \alpha & 1     & \alpha       & \alpha^2}.$$
Note that the symmetric power of the fixed basis $\set{c_1,c_2}$ is  $\set{c_{11},c_{12},c_{22}}$, which is a basis of $\Sym_2(M),$ so it is indeed rank $3$ as desired.

Consider the minor $B$ consisting of the following columns of $\Sym_2(M)$
$$B:=\kbordermatrix{
        &c_{11} & c_{22}    & c_{34}    \\
r_{11}  &1      & 0         & 0         \\
r_{12}  &0      & 0         & \alpha +1    \\
r_{22}  &0      & 1         & \alpha  }.$$
We can see that $B$ has determinant $-(\alpha+1)$. In particular if the characteristic is not $2$ then we may choose $\alpha= -1$ and $B$ will not be a basis. Otherwise, if we can choose $\alpha\neq 1,0,-1$ then $B$ will be a basis. Similarly, the matrix
$$C:=\kbordermatrix{
        &c_{11} & c_{22}    & c_{44}    \\
r_{11}  &1      & 0         & 1         \\
r_{12}  &0      & 0         & 2\alpha    \\
r_{22}  &0      & 1         & \alpha^2  }.$$
has determinant $-2\alpha$, which is 0 if and only if $k$ has characteristic 2. In particular is never a basis when represented in characteristic $2$ and it is always a basis otherwise.
\end{example}

\begin{remark}
\coloref{ex:matrixSymprodExample} shows that the underlying matroid of the second Macaulay matrix depends both on the field over which we take the representation $L$, as well as the representation $L$ itself over any fixed field. However, if $L$ is a representation of $M$ over the field $F$, then the second Macaualay matrix is preserved under action of $\GL_r(F)$ by left-multiplication in the following sense. If $U\in \GL_r(F)$ we have
    $$\Sym_2(UL) = \Sym_2(U)\Sym_2(L).$$
\coloref{ex:matrixSymprodExample} demonstrates that the analogue for right-hand multiplication by $V\in \GL_n(F)$ does not hold.
\end{remark}

We will prove in \coloref{cor:minorClosed} that the class of matroids with $d\th$ symmetric power is minor-closed. In light of \coloref{cor:tensorForbiddenMinors}, we can further show that the class of matroids with $d\th$ symmetric power has infinitely many forbidden minors
\begin{theorem}\label{thm:sympowForbiddenMinors}
    The class of matroids with $d\th$ symmetric powers has infinitely many forbidden minors.
\end{theorem}
\begin{proof}
    By \coloref{rem:truncation} it is sufficient to construct a family of forbidden minors for the class of matroids with second symmetric power. 

    Draisma and Rincon prove that if $M$ and $N$ have no second tensor power, then $M\oplus N$ has no second symmetric power \cite{DR19}. By \coloref{prop:VxU23}, for every $r\geq 4$, the matroids $M(r)\oplus U_{2,3}$ have no second symmetric power.

    We further claim that for integers $r_1>r_2\geq 4$, the matroids $M_1:=M(r_1)\oplus U_{2,3}$ and $M_2:=M(r_2)\oplus U_{2,3}$ have no common minors that are not realizable.
    
    We recall that, for all $r$, any contraction $M(r)\contr e$ is realizable, and all deletions common to $M(r_1)$ and $M(r_2)$ are realizable matroids \cite{Lind88}. In particular, a non-realizable minor of $M_1$ has circuits of size $r_1$, $r_1+1$, and at most one of size $3$ or less. Similarly, a non-realizable minor of $M_2$ has circuits of size $r_2$, $r_2+1$ and at most one of size $3$. Since $r_1>r_2>3$, these circuit sets cannot be equal and we conclude that there can be no common non-realizable minors between $M_1$ and $M_2$. The result follows.
\end{proof}

\subsection{Symmetric Powers by Flats}\label{subsec:sympowByFlats}
In \coloref{def:symprod} we decided on an axiomatization for matroid symmetric powers that is potentially stronger than that which is proposed by Mason in \cite{MGlueing}. Originally motivated from Lov\'asz's early work with symmetric powers of realizable matroids \cite{Lovasz77}, Mason proposed the following condition for a symmetric quasi power, as defined in  \coloref{def:symprod}, to be a symmetric power.
\begin{equation}\label{axiom:flatAxiom} 
\text{For every flat } F \text{ of } M \text{ the set } F\cdot \Sym_{d-1}(E) \text{ is a flat of } M_d.
\end{equation}

In \coloref{def:symprod} we instead required that \coloref{eqn:symprodRank} be satisfied.

\begin{proposition}\label{prop:SymprodImpliesFlats}
    If $M$ is a valuated matroid and $(M_1,\dots, M_d)$ is a $d\th$ symmetric power in the sense of \coloref{def:symprod}, then $M_d$ satisfies \hyperref[axiom:flatAxiom]{Axiom \ref*{axiom:flatAxiom}}.
\end{proposition}

We will prove \coloref{prop:SymprodImpliesFlats} in \coloref{subsec:sympowCharacterizations}, using our new geometric techniques to avoid an in-depth matroid theoretical argument. While we are able to prove that \coloref{def:symprod} implies \hyperref[axiom:flatAxiom]{Axiom \ref*{axiom:flatAxiom}}, the reverse is not at all obvious. Despite substantial effort, we were not able to make progress on this problem and have left it as an open problem left by Lov\'asz and Mason.

\begin{conjecture}(Lovasz '77, Mason '81)\label{conj:flatcondition}
If $(M_1,\dots,M_d)$ is a symmetric quasi power of $M$ and for every $F\in \sF(M)$ we have $F\cdot \Sym_{d-1}(E)\in \sF(M_d)$, then $M_d$ has rank $\binom{\rk(M)+d-1}{d}$.
\end{conjecture}

In order to prove the conjecture, it is sufficient to prove the following auxiliary result.

\begin{conjecture}\label{conj:flatcondition2}
    If $M$ is a matroid on the ground set $E$, $(M_1,\dots,M_d)$ is a symmetric quasi power of $M$, and for every $F\in \sF(M)$ we have $F\cdot \Sym_{d-1}(E)\in \sF(M_d)$, then there is an $e\in E$ such that the contraction
    $$M_d\contr (e\cdot \Sym_{d-1}(E))$$
    is a symmetric quasi power of $M\contr e$.
\end{conjecture}

\coloref{conj:flatcondition2} is nontrivial in general, as it is possible for $(M_1,\dots,M_d)$ to be a symetric quasi power, that is not full rank, and for the contraction $M_d\contr (e\cdot \Sym_{d-1}(E))$ to not be a symmetric quasi power for some $M\contr e$. However, it is not clear if this is necessarily possible that this could be the case for all $e$. It is not clear if \coloref{conj:flatcondition} is true, not even for second symmetric powers,  without additional substantial hypotheses.


\subsection{Further characterizations of Matroid Symmetric Powers}\label{subsec:sympowCharacterizations}

In this section we will discuss characterisations of symmetric powers, similar to how \coloref{lem:tensorCharacterizations} treats tensor products. In particular, we will provide a geometric characterization, alongside a characterisation by flats that requires a strictly stronger condition than that of \coloref{conj:flatcondition}.

In order to establish a geometric interpretation of matroid symmetric powers, we will need to define the \textit{variety} of a symmetric power.

\begin{definition}\label{def:matroidVariety}
Let $M$ be a matroid on ground set $[n]\cong\set{x_1,\dots,x_n}$, and let $M_d$ be a symmetric quasi power of $M$ on ground set $\Sym_{d}([n])\cong \Mon_d(n)$. For each circuit $C$ of $M_d$ we define the \textit{circuit polynomial}
    $$p_C := \bigoplus_{x\in C} x\in \T[x_1,\dots,x_n]_d.$$

We can now define the (tropical) \textbf{variety} of $M_d$ as
$$V(M_d) : = \bigcap_{C\in \sC(M_d)} V\left(p_C\right)\subset \T^n.$$
\end{definition}

Given a matroid $M$ on ground set $[n]$, its variety $V(M)\subset \T^n$ is equal to its \textbf{Bergman Fan}. For a symmetric power $M_d$ its variety and Bergman fan are distinct geometric objects. The Bergman fan treats each element of $\Sym_d(E)$ as an independent variable whereas the variety takes into account the algebraic dependencies between the variables. In particular, $B(M_d)\subset\T^{|\Sym_d(E)|}$ and $V(M_d)\subset \T^n$. We further remark that it is always the case that $V(M_d)\subset B(M)$ and that $V(M_d)$ maps to $B(M_d)$ through the Veronese embedding.


\begin{proposition}\label{prop:symprodCharacterisations}
Let $M_d$ be a $d\th$ symmetric quasi power of $M$. The following are equivalent for $M_d$
\begin{enumerate}[label=(\roman*),]
    \item\label{prop:symprodCharacterisations.rk} $M_d$ has rank $\binom{\rk(M)+d-1}{d}$
    \item\label{prop:symprodCharacterisations.bases} For every basis $B$ of $M$, $\Sym_d(B)$ is a basis of $M_d$
    \item\label{prop:symprodCharacterisations.bf} The variety $V(M_d)$ is supported on $B(M)$.
    \item\label{prop:symprodCharacterisations.flats} For every $\omega\in B(M)$ the set
    $$\set{x\in \Sym_d(E)\mid x(\omega)>0}$$
    is a flat of $M_d$
\end{enumerate}
\end{proposition}
The condition imposed by \coloritemref{prop:symprodCharacterisations}{flats} is stronger than \hyperref[axiom:flatAxiom]{Axiom \ref*{axiom:flatAxiom}}, which can be phrased in the following way.
\begin{quote}
    For every $\omega\in B(M)\cap \R^n_{\geq 0}$. The set
    $$\set{x\in \Sym_d(E)\mid x(\omega)>0}$$
    is a flat of $M_d$.
\end{quote}
Recall that the lineality space of $B(M)$ contains $\R \mathbbm{1}$. In particular, the positive part of the Bergman fan determines the entire fan. The discrepancy between \coloref{conj:flatcondition} and \coloritemref{prop:symprodCharacterisations}{flats} is a matter of whether or not we need to consider the flats corresponding to every point in the Bergman fan, or just the positive part --- the part that determines the entire fan.

\begin{proof} The equivalence between \coloritemref{prop:symprodCharacterisations}{rk}
and \coloritemref{prop:symprodCharacterisations}{bases} follows from \coloref{prop:masonFullRank}. We will prove the equivalence between \coloritemref{prop:symprodCharacterisations}{rk} and \coloritemref{prop:symprodCharacterisations}{bf} in \coloref{cor:bfCharacterisation}. We now prove the equivalence between \coloritemref{prop:symprodCharacterisations}{bf} and \coloritemref{prop:symprodCharacterisations}{flats}.

Suppose that $V(M_d)\neq B(M)$. Then $V(M_d)\subsetneq B(M)$ and so we let $\omega\in B(M)\sm V(M_d)$. Since $\omega\notin V(M_d)$ there is a circuit $C$ of $M_d$ that attains its minimum exactly once at $\omega$, say at the monomial $x$. Thus $x(\omega)$ is the unique minimum value of $C(\omega)$. Since the lineality space of $B(M)$ contains $\R\mathbbm{1}$ we may translate $\omega$ to $\omega' = \omega- \frac{x(\omega)}{d} \mathbbm{1}$ so that $x(\omega')=0$ is the unique minimum value of $C$ at $\omega'$. Now consider that this implies the set
$$S:=\set{y\in \Sym_d(E)\mid y(\omega')>0}$$
contains all of $C$ except for $x$. Thus $S$ is not a flat, as desired.

Now suppose that the set 
$$S = \set{y\in \Sym_d(E)\mid y(\omega)>0}$$
is a not a flat for some $\omega\in B(M)$. We claim that $V(M_d)\neq B(M)$. Since $S$ is not a flat there is a circuit $C$ such that $|C\sm S|=\set{x}$ --- a singleton. In particular, $x(\omega)$ is the unique minimum of $C$ at $\omega$ and so $\omega\notin V(M_d)$ and $V(M_d)\neq B(M)$, as desired.
\end{proof}

\begin{proof}[Proof of \coloref{prop:SymprodImpliesFlats}]
If $M$ is a matroid and $(M_1,\dots,M_d)$ is a symmetric power of $M$ as in \coloref{def:symprod}, then by \coloref{prop:symprodCharacterisations}, for every $\omega\in B(M)$ we have that

$$S_\omega:= \set{x\in \Sym_d(E): x(\omega)>0}.$$
Let $F$ be a flat of $M$. The set $F\cdot \Sym_{d-1}(E)$ is the set of elements of $\Sym_d(E)$ divisible by some element of $F$. Now let $\omega= e_F$. A monomial $x$ satisfies $x(\omega)>0$ if and only if $x$ is divisible by an element of $F$. In particular $F\cdot\Sym_{d-1}(E)=S_\omega$ and is therefore a flat of $M_d$, as desired.
\end{proof}

\subsection{Open Questions}\label{subsec:openQuestions}

We conclude this section by posing a few open questions that we were not successful in pursing in this paper. In the case of linear spaces there is a well-established connection between tensor and symmetric powers of linear spaces via the quotient of the action of the symmetric groups, we pose this as a question in the context of matroid powers

\begin{question}
Is there a relationship, e.g via matroid quotients, between matroid tensor powers and matroid symmetric powers? Are there suitable hypothesis on a matroid tensor power to guarantee the existence of a matroid symmetric power?
\end{question}

 We might further refine the previous question and ask if symmetric powers are, in some sense,  easier to construct than tensor powers.
\begin{question}\label{question:symprodsWithoutTensor}
Can a matroid have a $d \th$ symmetric power and no $d\th$ tensor power, or vice-versa?
\end{question}

A prime candidate for answering this question is to determine if the V\'amos matroid $V_8$ has a second symmetric power. We imagine that it does not, but if it does in fact have one, it would resolve this question and also prove that the class of matroids with second symmetric powers is not closed under duality. As a quick proof of this, if $V_8$ has a second symmetric power $P$, then so would $V_8\oplus U_{1,3}$ via $P\oplus (V_8\otimes U_{1,3})\oplus (U_{1,3})_2$, where $V_8\otimes U_{1,3}\cong V_8$. However, since $V_8$ is self-dual, the dual of $V_8\oplus U_{1,3}$ is $V_8\oplus U_{2,3}$, whose second symmetric power does not exist \cite{DR19}*{Theorem 5.2}. We refine this into its own question.

\begin{question}\label{question:vamosSym2}
Does the V\'amos matroid have a second symmetric power?
\end{question}
%%%%%%%%%%%%%%%%%%%%%%%%%%%%%%%%%%%%%%%%%%%%%%%%%%%%

%%%%%%%%%%%%%%%%%%%%%%%%%%%%%%%%%%%%%%%%%%%%%%%%%%%%

%%%%%%%%%%%%%%%%%%%%%%%%%%%%%%%%%%%%%%%%%%%%%%%%%%%%